\documentclass[a4paper,11pt,openany]{book}

\usepackage[margin=2.5cm]{geometry}
\usepackage[T1]{fontenc}
\usepackage{lmodern}

% More realistic TOC depth (chapters/sections/subsections)
\setcounter{tocdepth}{2}

% Simple page style
\pagestyle{plain}

% --- IMPORTANT: front/main/back matter WITHOUT resetting page counter ---
% book.cls uses \pagenumbering{roman/arabic} which resets page to 1.
% We redefine frontmatter/mainmatter/backmatter to only change \thepage (display),
% keeping the underlying page counter continuous.
\makeatletter
\renewcommand\frontmatter{%
  \if@openright\cleardoublepage\else\clearpage\fi
  \@mainmatterfalse
  \renewcommand\thepage{\roman{page}}%
}
\renewcommand\mainmatter{%
  \if@openright\cleardoublepage\else\clearpage\fi
  \@mainmattertrue
  \renewcommand\thepage{\arabic{page}}%
}
\renewcommand\backmatter{%
  \if@openright\cleardoublepage\else\clearpage\fi
  \@mainmatterfalse
  % keep arabic display in back matter (typical), and DO NOT reset counter
  \renewcommand\thepage{\arabic{page}}%
}
\makeatother

% ---- Filler text helpers (no extra packages needed) ----
\newcommand{\filler}{%
This is placeholder text for PDF reader tests. It is intentionally generic and repeated to
create realistic page flow, line breaks, paragraph boundaries, and searchable content.
It contains numbers (12345), punctuation, and mixed-case words for indexing tests. \par\medskip
}

% DO NOT use \count0 (it affects TeX's page output). Use our own counter:
\newcount\fillercount
\newcommand{\manyfiller}[1]{%
\begingroup
\fillercount=0
\loop\ifnum\fillercount<#1
  \advance\fillercount by 1
  \filler
\repeat
\endgroup
}

\begin{document}

% --------------------
% Front matter (roman numerals shown, counter NOT reset)
% --------------------
\frontmatter

% Cover page (anonymised)
\begin{titlepage}
  \centering
  \vspace*{2cm}
  {\Large \textbf{FREE eBook}\par}
  \vspace{2cm}
  {\Huge \textbf{LEARNING}\par}
  \vspace{0.2cm}
  {\Huge \textbf{Subject}\par}
  \vfill
  {\large Unaffiliated placeholder document\par}
  \vspace{0.5cm}
  {\large Version 1.0\par}
\end{titlepage}

\renewcommand{\contentsname}{Table of Contents}
\tableofcontents
\clearpage

\chapter*{About}
\addcontentsline{toc}{chapter}{About}
\noindent
About section (placeholder). This document is anonymised and intended only for PDF reader tests.
It includes headings, page breaks, lists, tables, code-like blocks, and footnotes.
\par\medskip
\filler
\filler

\chapter*{Credits}
\addcontentsline{toc}{chapter}{Credits}
Credits section (placeholder).\\
Contributor A, Contributor B, Contributor C.

% --------------------
% Main matter (arabic shown, counter continues)
% --------------------
\mainmatter

% ========== Chapter 1 ==========
\chapter{Chapter 1: Getting started}
\section{Overview}
\filler\manyfiller{3}

\section{Setup}
\subsection{Subsection 1}
\filler\manyfiller{2}

\subsection{Subsection 2}
\filler\manyfiller{2}

\section{Quick check}
A short list (placeholder):
\begin{itemize}
  \item Item 1 (placeholder)
  \item Item 2 (placeholder)
  \item Item 3 (placeholder)
\end{itemize}
\filler

% ========== Chapter 2 ==========
\chapter{Chapter 2: Basic concepts}
\section{Key terms}
\filler\manyfiller{2}

\subsection{Definition A}
\filler\manyfiller{2}

\subsection{Definition B}
\filler\manyfiller{2}

\section{A simple example}
Code-like block (placeholder):
\begin{verbatim}
# Example (placeholder)
command --flag value
output: OK
\end{verbatim}
\filler\manyfiller{2}

% ========== Chapter 3 ==========
\chapter{Chapter 3: Configuration}
\section{Settings}
\filler\manyfiller{2}

\subsection{Default values}
\filler\manyfiller{1}

\subsection{Overriding values}
\filler\manyfiller{2}

\section{A small table}
\begin{tabular}{|p{4cm}|p{9cm}|}
\hline
\textbf{Field} & \textbf{Description} \\
\hline
Field 1 & Placeholder description for field 1. \\
\hline
Field 2 & Placeholder description for field 2. \\
\hline
Field 3 & Placeholder description for field 3. \\
\hline
\end{tabular}

\par\medskip
\filler\manyfiller{2}

% ========== Chapter 4 ==========
\chapter{Chapter 4: Workflows}
\section{Common workflow}
\filler\manyfiller{2}

\subsection{Step 1}
\filler\manyfiller{2}

\subsection{Step 2}
\filler\manyfiller{2}

\section{Notes and remarks}
A quote (placeholder):
\begin{quote}
This is a placeholder quotation used to test paragraph indentation and block rendering.
\end{quote}
\filler\manyfiller{2}

% ========== Chapter 5 ==========
\chapter{Chapter 5: Data handling}
\section{Persistence}
\filler\manyfiller{3}

\subsection{Subsection 1}
\filler\manyfiller{2}

\subsection{Subsection 2}
\filler\manyfiller{2}

\section{Footnote test}
This sentence has a footnote\footnote{Footnote placeholder text for testing extraction and layout.}.
\filler\manyfiller{2}

% ========== Chapter 6 ==========
\chapter{Chapter 6: Networking}
\section{Connectivity modes}
\filler\manyfiller{3}

\subsection{Mode A}
\filler\manyfiller{2}

\subsection{Mode B}
\filler\manyfiller{2}

\section{Cross-reference text}
Cross-reference placeholder: see “Chapter 11: Troubleshooting”.
\filler

% ========== Chapter 7 ==========
\chapter{Chapter 7: Monitoring and logging}
\section{Monitoring}
\filler\manyfiller{3}

\section{Logging}
\filler\manyfiller{3}

\subsection{Example log snippet}
\begin{verbatim}
[INFO] Placeholder log line 1
[WARN] Placeholder log line 2
[ERROR] Placeholder log line 3
\end{verbatim}
\filler\manyfiller{2}

% ========== Chapter 8 ==========
\chapter{Chapter 8: Security}
\section{Principles}
\filler\manyfiller{3}

\subsection{Checklist}
\begin{enumerate}
  \item Checklist item 1 (placeholder)
  \item Checklist item 2 (placeholder)
  \item Checklist item 3 (placeholder)
\end{enumerate}
\filler\manyfiller{2}

% ========== Chapter 9 ==========
\chapter{Chapter 9: Advanced topic}
\section{Deep dive}
\filler\manyfiller{6}

\subsection{Subtopic}
\filler\manyfiller{3}

% ========== Chapter 10 ==========
\chapter{Chapter 10: Reference}
\section{Commands}
\filler\manyfiller{2}

\subsection{Command A}
\begin{verbatim}
command-a --help
\end{verbatim}
\filler\manyfiller{2}

\subsection{Command B}
\begin{verbatim}
command-b --verbose --dry-run
\end{verbatim}
\filler\manyfiller{2}

% ========== Chapter 11 ==========
\chapter{Chapter 11: Troubleshooting}
\section{Build issues}
\filler\manyfiller{3}

\subsection{Symptom}
\filler\manyfiller{2}

\subsection{Resolution}
\filler\manyfiller{3}

\section{Runtime issues}
\filler\manyfiller{4}

% --------------------
% Back matter (arabic continues)
% --------------------
\backmatter

\chapter*{Appendix A: Extra notes}
\addcontentsline{toc}{chapter}{Appendix A: Extra notes}
\filler\manyfiller{4}

\chapter*{Appendix B: Glossary}
\addcontentsline{toc}{chapter}{Appendix B: Glossary}
\begin{description}
  \item[Term 1] Definition (placeholder).
  \item[Term 2] Definition (placeholder).
  \item[Term 3] Definition (placeholder).
\end{description}
\filler\manyfiller{2}

\chapter*{Index (placeholder)}
\addcontentsline{toc}{chapter}{Index (placeholder)}
Index placeholder text.

\end{document}